% TOP_PROBLEM: {"id": 30006958, "title": "Polynomial-Time Quasioptimal Logarithmic-Energy Points on the 2-Sphere", "category_id": 20, "category": {"id": 20, "name": "miscellaneous", "display_name": "Miscellaneous", "description": "Problems whose source classification does not fit the main mathematical categories.", "slug": "miscellaneous", "order_index": 20, "created_at": "2026-07-31T15:26:25.671Z"}, "status": "open"}
% FIRST_POSTED: 2026-09-27
% !TeX program = lualatex
% ATTEMPT_STATUS: unresolved
\documentclass[11pt]{article}
\usepackage[margin=25mm]{geometry}
\usepackage{fontspec}
\setmainfont{Latin Modern Roman}
\usepackage{amsmath,amssymb,mathtools}
\usepackage{unicode-math}
\setmathfont{Latin Modern Math}
\usepackage{xurl,hyperref}
\hypersetup{colorlinks=true,urlcolor=blue}
\setcounter{secnumdepth}{0}
\setlength{\parindent}{0pt}
\setlength{\parskip}{5pt}
\setlength{\emergencystretch}{3em}
\begin{document}
\section*{\texorpdfstring{Polynomial-Time Quasioptimal Logarithmic-Energy Points on the 2-Sphere}{Polynomial-Time Quasioptimal Logarithmic-Energy Points on the 2-Sphere}}\label{30006958}
\subsection{Definitions and mathematical statement}
For distinct points $x_1,\ldots,x_N$ on the unit sphere $S^2\subset{\mathbb{R}}^3$, define logarithmic energy, using Smale's pair-counting convention, by
\[
 E(x_1,\ldots,x_N)=\sum_{1\le i<j\le N}
 \log\frac1{\|x_i-x_j\|},\qquad E_N^{\min}=\inf E.
\]
Seek a uniform real-number algorithm and an absolute constant $c$ such that, for every $N\ge2$, it produces $N$ points in time polynomial in $N$ and
\[
 E(x_1,\ldots,x_N)\le E_N^{\min}+c\log N.
\]
The real-number computation model and admissible operations are those in Smale's Problem 7. This does not silently impose a bit-complexity bound for exact transcendental coordinates. The required error is additive and logarithmic, much smaller than a fixed relative error in the leading-order energy.
\par\smallskip\textit{Formulation and concept references:} \href{https://www.cityu.edu.hk/ma/doc/people/smales/pap104.pdf}{[S1]}.

\subsection{Short English statement}
Can a polynomial-time real-number algorithm place points on the sphere with logarithmic energy within a constant times log N of optimal?
\subsection{Research attempt}
\paragraph{Scope and source audit (27 September 2026).}
This attempt keeps the unordered-pair convention in the statement. A
paper summing over $i\ne j$ uses twice this energy, including twice every
additive error. The requested algorithm is uniform in $N$, and its error
must be compared with the global minimum, not with the best configuration
found in a numerical run. Polynomial time for a local optimization step
does not bound the number of steps or certify the global energy gap.

The original hosted PDF in [S1] was not retrievable during this audit.
The primary paper [E1], Problem 1.1, explicitly restates the logarithmic
additive target and develops lower bounds; [E2] computes expected energies
for a family of random constructions. Neither inspected result supplies
the requested algorithm. The discussion below does not treat a random
expectation or an asymptotic leading coefficient as that guarantee. It
also does not impose exact transcendental output in a bit model that is
absent from the catalogue statement.

The approach is deliberately elementary: derive exact small cases and a
global separation estimate, then discretize the continuous optimization.
The discretization reaches a constant additive error using rational
coordinates of moderate description length. Its search has exponentially
many candidates. This isolates a selection problem rather than an
obstruction from coordinate precision.

\paragraph{Existence of a minimizer and the product formulation.}
Put $m=\binom N2$. Every pair distance is at most two, so every summand
is at least $-\log2$. There is at least one distinct configuration of
finite energy. Along a sequence in which any two points collide, the
corresponding summand tends to $+\infty$, while the other summands have
the fixed lower bound just noted. Thus a bounded-energy minimizing
sequence cannot converge to a collision. Compactness of $(S^2)^N$ and
continuity away from collisions prove that $E_N^{\min}$ is attained.

Equivalently, maximize
\begin{equation}\label{a474:product}
 D(X)=\prod_{i<j}\|x_i-x_j\|^2,\qquad E(X)=-\tfrac12\log D(X).
\end{equation}
For rational sphere points, $D(X)$ is a positive rational number.
Comparing two such energies therefore requires only exact rational
multiplication and comparison, not exact evaluation of logarithms.
The requested gap is equivalent to
\[
 D(X)\ge D(X_{\min})N^{-2c}.
\]
This is a multiplicative approximation to the entire distance product;
approximating each individual distance to a fixed relative tolerance
would accumulate an error over $m$ pairs.

\paragraph{Lemma 474.1: a centroid bound and three exact optima.}
For every $N$-point configuration,
\begin{equation}\label{a474:jensen}
 E(X)\ge-\frac m2\log\frac{2N}{N-1}.
\end{equation}
Equality holds precisely when the centroid is zero and all pair distances
are equal. Consequently the bound is attained for $N=2,3,4$ by the
antipodal pair, equilateral equatorial triangle, and regular tetrahedron.

\emph{Proof.} Expanding scalar products gives
\[
 \sum_{i<j}\|x_i-x_j\|^2
 =N\sum_i\|x_i\|^2-\left\|\sum_i x_i\right\|^2
 =N^2-\left\|\sum_i x_i\right\|^2\le N^2.
\]
Concavity of the logarithm yields
\[
 \frac1m\sum_{i<j}\log\|x_i-x_j\|^2
 \le\log\left(\frac1m\sum_{i<j}\|x_i-x_j\|^2\right)
 \le\log\frac{2N}{N-1}.
\]
Strict concavity gives the stated equality conditions. If the points
form an equal-distance, zero-centroid configuration, their Gram matrix
has diagonal one and off-diagonal $-1/(N-1)$. Its rank is $N-1$;
because the vectors lie in $\mathbb R^3$, this requires $N\le4$.
The three listed configurations have exactly that Gram matrix.
\hfill$\square$

In particular
\[
 E_2^{\min}=-\log2,\qquad E_3^{\min}=-\tfrac32\log3,
 \qquad E_4^{\min}=-3\log(8/3).
\]
The rank restriction explains why simply imposing the equality conditions
does not solve larger instances. Zero centroid alone is far weaker than
equality: many very poorly distributed configurations have zero centroid.

\paragraph{An exact average potential, with its algorithmic limitation.}
Let $\sigma$ be normalized surface area measure. Rotational invariance
makes $I=\int_{S^2}\log(1/\|x-y\|)\,d\sigma(x)$ independent of $y$.
Taking $y$ to be the north pole, the height $u=x\cdot y$ is uniform on
$[-1,1]$. Direct integration, including the integrable endpoint singularity,
gives
\begin{equation}\label{a474:average}
 I=-\frac14\int_{-1}^1\log(2-2u)\,du
   =\frac12-\log2<0.
\end{equation}
For independent uniform points, collisions have probability zero and
linearity of expectation gives $\mathbb E E=mI$. Therefore
\begin{equation}\label{a474:upper}
 E_N^{\min}\le mI.
\end{equation}

There is also an existence-only greedy version. Given $n-1$ distinct
points, the potential $V(x)=\sum_{j<n}\log(1/\|x-x_j\|)$ has integral
$(n-1)I$, so some noncolliding point satisfies $V(x)\le(n-1)I$.
Choosing such a point at each stage yields energy at most $mI$.
This observation neither tells us how to find that point efficiently in
the allowed real-number model nor narrows the error to $O(\log N)$.
The proven upper and lower bounds themselves leave an order-$N^2$ interval.

A useful failed construction is the regular equatorial $N$-gon. If
$z_1,\ldots,z_N$ are the roots of $z^N-1$, differentiation at each root
gives $\prod_{j\ne i}|z_i-z_j|=N$. Multiplying over $i$ shows
$\prod_{i<j}|z_i-z_j|^2=N^N$, hence
\[
 E(\text{regular equatorial }N\text{-gon})=-\tfrac N2\log N.
\]
Combining with \eqref{a474:upper}, its gap is at least
$-\tfrac N2\log N-mI$, which grows quadratically. Exact symmetry and
an exact distance product do not imply near-optimality on the sphere.

\paragraph{Lemma 474.2: an elementary separation bound at a minimizer.}
Every minimizing configuration has pairwise separation at least
\begin{equation}\label{a474:separation}
 \delta_N=2\exp(-(N-1)/2).
\end{equation}
This bound is intentionally weak, but sufficient for the discretization.

\emph{Proof.} Fix point $x_i$ of a minimizer. Replacing just this point
cannot lower the total energy, so its potential is at most the spherical
average of that same potential:
\[
 \sum_{j\ne i}\log(1/\|x_i-x_j\|)\le(N-1)I.
\]
Equivalently $\prod_{j\ne i}\|x_i-x_j\|\ge e^{-(N-1)I}$.
All but any selected factor are at most two. Thus, for every $j\ne i$,
\[
 \|x_i-x_j\|\ge 2^{-(N-2)}e^{-(N-1)I}
                  =2e^{-(N-1)/2}.
\]
The averaging argument is legitimate despite the finitely many infinite
values at the other points: the potential is integrable and its global
minimum lies away from those singularities.\hfill$\square$

No claim is made that this is the sharp scale. In particular, using this
exponentially small bound in a covering construction will produce a very
large grid. A polynomial separation estimate would improve the grid size,
but would still leave the selection of $N$ points from that grid.

\paragraph{Lemma 474.3: perturbation cost at a separated configuration.}
Suppose $X=(x_i)$ has separation at least $\delta>0$, and
$Y=(y_i)$ satisfies $\|x_i-y_i\|\le\varepsilon$ for every $i$, with
$2\varepsilon\le\delta/2$. Then the $y_i$ are distinct and
\begin{equation}\label{a474:perturb}
 E(Y)-E(X)\le4m\varepsilon/\delta.
\end{equation}

\emph{Proof.} The reverse triangle inequality gives
$\|y_i-y_j\|\ge\|x_i-x_j\|-2\varepsilon>0$.
The increase for that pair is at most
$-\log(1-2\varepsilon/\|x_i-x_j\|)$.
For $0\le u\le1/2$, integration of $1/(1-u)\le2$ gives
$-\log(1-u)\le2u$. Sum over the $m$ pairs.\hfill$\square$

Taking $\varepsilon\le\delta_N/(4N^2)$ makes the increase less than
$1/2$. Thus a sufficiently fine rational approximation of a minimizing
configuration already meets a stronger energy tolerance than requested.
The issue is obtaining such an approximation without knowing a minimizer.

\paragraph{A finite rational search that attains a constant gap.}
Here are explicit charts, avoiding an unmentioned transcendental-coordinate
oracle. The rational parametrization
\[
 F(u,v)=\frac{(2u,2v,1-u^2-v^2)}{1+u^2+v^2}
\]
maps $\mathbb R^2$ into $S^2$. On $[-1,1]^2$ its differential has
operator norm at most two. Use this chart with either pole and permute
the three coordinate axes. These six charts cover the sphere: choose a
coordinate of maximal absolute value and use the corresponding hemisphere;
the inverse stereographic coordinates then have absolute values at most one.

Let $K=4N+10$ and take the grid $u,v\in2^{-K}\mathbb Z\cap[-1,1]$ in
each chart. Rounding chart coordinates gives a covering radius at most
$2\sqrt2\,2^{-K}$. For $N\ge2$ this is less than
$\delta_N/(4N^2)$. One can check the inequality without an exponential
oracle by using $e<4$ in \eqref{a474:separation}; the deliberately generous
choice of $K$ leaves ample slack. Every grid point has rational coordinates
with $O(N)$ bits per numerator and denominator.

Enumerate all distinct $N$-subsets of the resulting finite point set and
select one maximizing \eqref{a474:product}. The grid contains a rounding
of some minimizing configuration. Lemma 474.3 shows that its energy is
less than $E_N^{\min}+1/2$, so the selected subset has at least this
quality. In particular this proves existence of rational, moderate-length
descriptions achieving the desired energy scale.

There are $2^{O(N)}$ grid points and at most $2^{O(N^2)}$ ordered candidate
tuples. Exact products contain $O(N^3)$ bits, up to fixed constants,
since there are $O(N^2)$ factors with $O(N)$-bit coordinates. Exhaustive
search therefore gives an exponential-time construction with polynomial
size output; it is not the requested polynomial-time real-number algorithm.
The same conclusion holds if arithmetic operations themselves have unit
cost: the number of subsets visited is still exponential.

\paragraph{Why the natural shortcuts do not remove the search.}
The energy is not a convex function of the sphere-point coordinates on
their nonconvex domain. A stationary configuration or a descent direction
does not compare its value with the global minimum. Nor does the compactness
argument give an efficient covering of the set of minimizers. The rational
grid construction supplies a verifiable objective value, but its proof
of quality uses the exhaustive maximum over all subsets.

One might instead compare a construction with an asymptotic expansion of
$E_N^{\min}$. An error term $O(N)$, or even an unspecified $o(N)$ term,
does not imply $O(\log N)$: for example $N^{1/2}=o(N)$ exceeds every
constant multiple of $\log N$. Likewise the word ``quasioptimal'' in a
different source may mean a bounded ratio or a linear additive error.
The exact inequality in this catalogue is the criterion used throughout.

A sufficient missing lemma would select, in polynomial time, an $N$-subset
of an explicitly generated sphere grid whose distance product is within a
polynomial factor of the best continuous product. A grid with a polynomial
number of points would not alone prove that lemma: optimizing over its
$N$-subsets can still involve exponentially many choices. Neither such a
selection theorem nor a rigorous global convergence bound is established.

\paragraph{Finite diagnostics and outcome.}
The companion checks use exact rational sphere points from the charts,
verify their norms and the distance-product comparison, and compare
enumerated small-grid optima with the exact two-, three-, and four-point
bounds. The tetrahedral Gram calculation is checked independently from
its coordinate realization. The finite tests do not estimate an unknown
large-$N$ optimum or certify a numerical local optimizer.

\textbf{Outcome: UNRESOLVED.} The attempt proves exact small subcases,
a global separation estimate, a perturbation bound, and an explicit
constant-gap rational search with exponential running time. It pinpoints
efficient global selection as the first remaining algorithmic gap.
Nothing here supplies the polynomial-time logarithmic-gap algorithm.


\subsection{Sources}
\begin{itemize}
\item[S1]Stephen Smale, Mathematical Problems for the Next Century, Problem 7 (1998).
\url{https://www.cityu.edu.hk/ma/doc/people/smales/pap104.pdf}.
\textit{Formulation source.}
\item[E1]Carlos Beltr\'an and F\'atima Lizarte,
\emph{A Lower Bound for the Logarithmic Energy on $S^2$ and for the Green
Energy on $S^n$}, Constructive Approximation, published online 2023,
Problem 1.1 and introduction.
\url{https://link.springer.com/article/10.1007/s00365-023-09642-4}.
\textit{Primary formulation and lower-bound comparison inspected
27 September 2026; the unavailable original link [S1] is preserved.}
\item[E2]\emph{On the logarithmic energy of solutions to the polynomial
eigenvalue problem}, arXiv:2408.11148 (2024).
\url{https://arxiv.org/abs/2408.11148}.
\textit{Primary abstract inspected 27 September 2026; used only for the
scope of random expected-energy constructions, not a worst-case algorithm.}

\end{itemize}
\end{document}
